\documentclass[12pt]{article}
\usepackage[utf8]{inputenc}

\usepackage{mystyle}
\usepackage{parskip}
\usepackage{float}
\usepackage[normalem]{ulem}

% Comic Sans
% \usepackage{sansmathfonts}
% \usepackage{fontspec}
% \setmainfont{Comic Sans MS}

% \usepackage{fontspec}
% \setmainfont{Wingdings-Regular.ttf}

% \usepackage{libertine}

% fancy header
\usepackage{fancyhdr}
\pagestyle{fancy}
\rhead{Joseph Sullivan}

\title{Projective Space}
\author{Joseph Sullivan}
\date{August 2024}



\begin{document}

\maketitle

\tableofcontents

\begin{itemize}
    \item sheaves on $\P^n$ vs graded module
    \begin{itemize}
        \item correspondence up to sufficiently high degree
        \item resolutions
    \end{itemize}
\end{itemize}

\section{Proj Construction}

Let $S$ be a ($\Z^{\geq 0}$-)graded ring.

\begin{rmk}
    If $T\subseteq S$ is a multiplicative set consiting of homogeneous elements, then $T^{-1}S$ is naturally a $\Z$-graded ring, where we set
    \[\deg \frac{g}{t} = \deg g - \deg t.\]
\end{rmk}

In analogy to the classical setting, where $S=\C[x_0,\dots,x_n]/I$ for $I$ a homogeneous ideal, we will define $\Proj S$ by gluing together the dehomogenizations of $S$ along coordinates. To generalize, however, we allow \textit{any} coordinate, not just the ``standard'' coordinates $x_i$.

Let's recall how dehomogenization works classically. For example, for example, consider $S=\C[x_0,x_1,x_2]/(g)$, where
\[g(x_0,x_1,x_2)=x_0 x_1 + x_2^2.\]
Classically, we pick a coordinate, say $x_0$, and look at the defining equation of $g$ in the complement of its vanishing $U_0 = \P^3\setminus V(x_0)$. We would realize that in homogeneous coordinates we can rescale the nonzero $x_0$ to be $1$, which is effectively dividing the coordinates by $x_0$. So, we would evaluate
\[g(\frac{x_0}{x_0}, \frac{x_1}{x_0}, \frac{x_2}{x_0}) = \frac{x_0}{x_0}\frac{x_1}{x_0} + \left(\frac{x_2}{x_0}\right)^2,\]
and we might rename these new variables $x_{1/0}=\frac{x_1}{x_0}$, $x_{2/0}=\frac{x_2}{x_0}$. Then, our chart with respect to $x_0$ is
\[\Spec \C[x_{1/0},x_{2/0}]/(x_{1/0}+x_{2/0}^2).\]
This coordinate ring is exactly the degree 0 part of the ring $(\C[x_0,x_1,x_2]/g)_{x_0}$.

Now, how should this work more generally? We want a notion of dehomogenizing at an arbitrary homogeneous $f\in S_+$, which we interpret as a coordinate. Now, we again study the complement of the vanishing of the coordinate, which is captured algebraically by $S_f$. Then, we look at the degree 0 part---this is because we want a standard representative of functions which are all equivalent. So, we will declare
\[\Spec (S_f)_0\]
to be a chart of $\Proj S$ (which is still yet to be defined).

\textbf{As a set}, define
\[\Proj S = \{\text{homogeneous prime ideals of $S$ not containing $S_+$}\}.\]

\begin{prop}
    Let $A$ be a $\Z$-graded ring with an invertible positive degree element $f$ (e.g., $A=S_f$). Then, there is a bijection between the prime ideals of $A_0$ and the homogeneous prime ideals of $A_f$ by
    \begin{align*}
        \{\text{homog. primes of $A$}\} &\leftrightarrow \{\text{primes of $A_0$}\}\\
        P &\mapsto P\cap A_0\\
        \bigoplus Q_i &\mapsfrom P_0
    \end{align*}
    where $Q_i\subseteq A_i$ is the set
    \[Q_i = \{a \in A_i \mid a^{\deg f}/f^i \in P_0\}.\]
\end{prop}

\begin{proof}
    The left-to-right direction gives a well-defined prime ideal as usual for a map of rings $A_0 \hookrightarrow A$.

    We show that the right-to-left is well-defined, so we show
    \[Q = \bigoplus Q_i\]
    is a homogeneous prime ideal. Once we show it is an ideal, it is automatically homogeneous by its construction (generated by homogeneous forms in each $Q_i$).

    First, we prove a lemma, that $a\in Q_i$ iff $a^2\in Q_{2i}$. If $a\in Q_i$, then $a^{\deg f}/f^i \in P_0$, and since $P_0$ is an ideal, we can scale it to get
    \[\frac{a^{\deg f}}{f^i} \cdot \frac{a^{\deg f}}{f^i} = \frac{(a^2)^{\deg f}}{f^{2i}} \in P_0,\]
    so $a^2 \in Q_{2i}$. Conversely, if $a^2 \in Q_{2i}$, then $P_0$ being prime and the previous equation gives $a \in Q_i$.
    
    Now, to show $Q$ is at least an abelian subgroup, we just need to show each $Q_i$ is an abelian subgroup, so let $a,b\in Q_i$. To show $a+b\in Q_i$, we show $(a+b)^2 \in Q_{2i}$. We have
    \[\frac{((a+b)^2)^{\deg f}}{f^{2i}} = \frac{(a + b)^{2\deg f}}{f^{2i}},\]
    and when we distribute, by pigeonhole, each term will have at least one of $a$ or $b$ a power at least $\deg f$ (without loss of generality, say $a$), so
    \[\frac{a^m b^n}{f^{2i}} = \frac{a^{m-\deg f} b^n}{f^i} \frac{a^{\deg f}}{f^i},\]
    which is in $P_0$ because it is an ideal, so we get that $a+b\in Q_i$. Thus, $Q_i$ is indeed an abelian group, and it is not hard to show that it is an ideal, so the right-to-left construction is well-defined.

    Now, we show that this actually gives a bijection, i.e., that the directions are inverse to each other. For the composition
    \[P_0 \mapsto \bigoplus Q_i \mapsto Q_0,\]
    by definition $Q_0 = \{a\in A_0 \mid a^{\deg f} \in P_0\}$, which tells us
    \[P_0 \subseteq Q_0 \subseteq \sqrt{P_0},\]
    but $P_0$ is radical anyways, so $Q_0 = P_0$.

    Conversely, for the composition
    \[P \mapsto P\cap A_0 \mapsto \bigoplus Q_i,\]
    we will at least have $\bigoplus Q_i \subseteq P$ since if $a^{\deg f}/f^i \in P \cap A_0$, then $a^{\deg f} \in P$, so by primeness $a\in P$. Then, for $a\in P$ is homogeneous, we have $a^{\deg f}/f^{\deg a} \in P\cap A_0$, so $a\in Q_{\deg a}$, which then implies $P\subseteq \bigoplus Q_i$.

    Thus, these maps are indeed inverse, and the proposition holds.
\end{proof}

This tells us that our definition of $\Proj S$ as a set respects dehomogenization, and that we should/will be able to identify $\Spec (S_f)_0$ with a subset of $\Proj S$.

\textbf{As a topological space}, define the closed sets of $\Proj S$ to be, for $I\leq S$ a homogeneous ideal,
\[V(I) = \{P \in \Proj S \mid P \geq I\},\]
and one should also check that this makes $\Spec (S_f)_0 \to \Proj S$ an inclusion of an open set. In fact, defining
\[D(f) = \Proj S \setminus V(f),\]
one should see that $\Spec (S_f)_0$ factors through $D(f)$, and that this is a homeomorphism.

\textbf{As a ringed space}, define $\cO_{\Proj S}$ by gluing $\cO_{D(f)} = \cO_{\Spec (S_f)_0}$ along intersections $D(fg)$ (checking the cocycle condition on $D(fgh)$).

\subsection{Homogeneous Tilde}

Let $M$ be a $\Z$-graded $S$-module. Define the quasicoherent sheaf $\tilde{M}$ on $\Proj S$ by defining it on the base of open sets $D(f)$ for $f\in S$ homogeneous. Set
\[\tilde{M}(D(f)) = (M_f)_0.\]
In particular, a $\Z^{\geq 0}$-graded ring $S$ comes equipped with the graded module $S^+$. When $S$ is standard graded (i.e., generated in degree 1), we will denote
\[\cO_{\Proj S}(1) \defeq \tilde{S_+}.\]
\begin{prop}
    $\cO_{\Proj S}(1)$ is a line bundle.
\end{prop}
\begin{proof}
    We have to trivialize $\cO_{\Proj S}(1)$. Since $S$ is generated in degree 1, $S^+ = \sum_{i\in I} S\cdot s_i$ for $s_i\in S_1$ and $I$ some index set. This tells us that 
    \[\Proj S = \bigcup_{i\in I} D(s_i),\]
    since any homogeneous prime ideal not containing $S^+$ must miss some $s_i$. Then, on $D(s_i)$, the sheaf $\cO_{\Proj S}(1)$ has a no where vanishing section $s_i$, which gives an isomorphism
    \[\cO_{D(s_i)} \xrightarrow{\cdot s_i} \cO_{\Proj S}(1)|_{D(s_i)},\]
    so indeed $\cO_{\Proj S}(1)$ is a line bundle.
\end{proof}

\section{Morphisms to Projective Space}

\subsection{Some Basics}

\subsubsection{Using the Trivial Line Bundle}

As a first attempt to give maps to projective space, let $X$ be a $k$-scheme, and $f_0,\dots,f_r\in \cO_X(X)$ some global functions that don't simultaneously vanish. Then, we want to define a map
\[X \xrightarrow{[f_0: \cdots : f_r]} \P^r_k.\]
At least at $k$-points, we get a map of sets, where we send $x$ to the (homogeneous maximal ideal corresponding to) $k^*$ equivalence class $[f_0(x),\dots,f_r(x)]$, where $f_i(x)\in k$. To make more sense out of this map, we should say what happens to other points, and what happens to the structure sheaf. To do this, we locally define
\[D(f_i) \xrightarrow{(\frac{f_0}{f_i}, \dots, \hat{\frac{f_i}{f_i}}, \dots, \frac{f_r}{f_i})} D(x_i) \subseteq \P^r_k,\]
where the $D(f_i)$ cover $X$ because $V(f_0,\dots,f_r)=\emptyset$. Now, this morphism makes sense (on the level of schemes) by the $\Gamma, \Spec$ adjunction, where we are giving a ring homomorphism
\[\Gamma(D(x_i)) = (k[x_0,\dots,x_r]_{x_i})_0 \cong k[x_{0/i},\dots,\hat{x_{i/i}},\dots,x_{r/i}] \longrightarrow \Gamma(D(f_i))\]
by sending $x_{j/i} \mapsto f_j/f_i \in \Gamma(D(f_i))$.

However, we don't have much hope of producing many morphisms in this way. For example, we cannot even produce the morphism $\P^r_k \xrightarrow{\id} \P^r_k$, since $\P^r_k$ has only $k$ for global sections. To have better luck, we need to use sections not just from $\cO_X$, but from more complicated line bundles.

\subsubsection{Basepoint-free Line Bundles}

First, we discuss some adjectives about sheaves and their sections. Let $X$ be a $k$-scheme and $\cF$ a sheaf. We say $V\subseteq H^0(X,\cF)$ \textbf{generates} $\cF$ if
\[V \otimes_k \cO_X \xrightarrow{\eval} \cF\]
is surjective. If some $V$ generates $\cF$ (in particular $V=H^0(X,\cF)$ all global sections), we say $\cF$ is \textbf{globally generated}.

In a similar vein, we say $\cF$ is \textbf{globally generated at $P$} by $V$ if evaluation of $V$ at the stalk of $P$ is surjective, i.e.,
\[V \otimes_k \cO_{X,P} \twoheadrightarrow \cF_p.\]

\begin{prop}
    Let $\cF$ be a finite type (i.e., sections over affine open is finite generated) quasicoherent sheaf. Then, $\cF$ is globally generated at $P$ by $V$ iff global sections generate the fibers at $P$ by $V$, i.e.,
    \[V \otimes_k \cO_{X,P}/\fm_P\cO_{X,P} \to \cF_P/\fm_P\cF_P\]
    is surjective.
\end{prop}
\begin{proof}
    If $V \otimes_k \cO_{X,P} \twoheadrightarrow \cF_P$, then restrictions of sections in $V$ to $\cF_P$ generate the stalk, so their images generate the quotient $\cF_P/\fm_P \cF_P$. Alternatively, we could view this as tensoring by the skyscraper sheaf $k(P)$ is right exact, so preserves surjections.

    Conversely, if $V \otimes_k \cO_{X,P}/\fm_P \cO_{X,P} \twoheadrightarrow \cF_P/\fm_P \cF_P$, then there are global sections $m_1,\dots,m_n\in V$ of $\cF$ which restrict/quotient to generate the vector space $\cF_P/\fm_P \cF_P$, the residue field of a finitely generated $\cO_P$-module $\cF_P$. So, by Nakayama's lemma, $m_1,\dots,m_n$ generate $\cF_P$.
\end{proof}

\begin{prop}
    If $\cF$ is finite type and is globally generated at $P$ by $V$, then there is a neighborhood $U$ of $P$ where $\cF$ is globally generated at all points of $U$.
\end{prop}

\begin{proof}
    Since $\cF$ is finite type and is globally generated at $P$ by $V$, there are sections $s_1,\dots,s_r \in V$ with $\<s_1,\dots,s_r\>=\cF_P$. We want to show that $s_1,\dots,s_r$ generate the fibers for nearby $Q$ (using the previous proposition).

    In an affine open $\Spec A \subseteq X$, $\cF =\tilde{M}$ for a finitely generated $A$-module. We have $M_P$ generated by the images of elements $m_1,\dots,m_r\in M$ associated to $s_1,\dots,s_r$, and we want to show that they generate nearby $M_Q$. 
\end{proof}

In particular, if $\cF = L$ is a line bundle, then $L_P/\fm_P L_P$ is a 1-dimensional $k(P)$-vector space, so this lemma specializes to the statement that $L$ is globally generated at $P$ by $V$ iff there exists a section $f\in V$ which has $f(P)\neq 0$. When this fails, we say that $P$ is a \textbf{base point} for $V$.

\begin{prop}
    Let $X$ be quasicompact. If $\cF$ is globally generated by $V\leq H^0(X,\cF)$ at all closed points of $X$, then $\cF$ is globally generated by $V$ at all points of $X$.
\end{prop}

\begin{proof}
    This is a standard specialization to closed points type argument. Let $P\in X$, and since $X$ is quasicompact, there is a closed point specialization $\fm\in \overline{\{P\}}$.

    Since $\cF$ is globally generated by $V$ at closed points, we have
    \[\cO_{X,\fm} \otimes_k V \twoheadrightarrow \cF_\fm,\]
    and localizing at the prime ideal $P\cO_{X,\fm}$, we get
    \[\cO_{X,P} \otimes_k V \twoheadrightarrow \cF_P,\]
    since localization is exact.

    Heuristically, this proof is saying that some global sections which, doing $\cF_\fm$ linear combinations, globally generate $\cF$ at $\fm$ will also globally $\cF$ at $P$ (the generic point of a subscheme), where we do $\cF_P$ linear combinations, which are much more plentiful (generating much more).
\end{proof}

\subsubsection{Embedding using a Line Bundle}
Now, suppose $V$ generates $L$. We show how this gives rise to a morphism from $X$ to projective space.

\textbf{Non-intrinsic way.} First, we do it in a choice-y way, until we later show our choice doesn't really matter. Choose a basis $V=\<s_0,\dots,s_r\>$. In analogy to our discussion about morphisms using the trivial line bundle, we will define
\[X \xrightarrow{[s_0: \cdots : s_r]} \P^r_k.\]
As before, this makes sense on $k$-points, since the tuple $(s_0(P),\dots,s_r(P)) \in (L_P/\fm_P L_P)^{r+1}$ can be identified (up to $k^*$ action) with a nonzero tuple in $k^{r+1}$, and then we can pass to its equivalence class.

But, as before, to make full sense of this we will work locally. Since $V$ is basepoint free (because it generates $L$), the open sets $D(s_i)$ cover all of $X$. On each $D(s_i)$, the section $s_i$ is nonvanishing, so
\[\cO_{D(s_i)} \xrightarrow{\cdot s_i} L|_{D(s_i)}\]
is an isomorphism (it is surjective on fibers, so on stalks, so surjective, and a surjective map of line bundles is an isomorphism).

We will abuse notation and denote the inverse of this morphism of sheaves by $s\mapsto \frac{s}{s_i}$. Now, we can define
\[D(s_i) \xrightarrow{(\frac{s_0}{s_i}, \cdots, \hat{\frac{s_i}{s_i}}, \cdots, \frac{s_r}{s_i})} D(x_i) \subseteq \P^r_k.\]
One should then check that this is well-defined on all of $X$ (that these morphisms restrict to the same morphism on pairwise intersections $D(s_i s_j)$).

\textbf{Intrinsic way.} Now, we show an intrinsic way that $V$ gives a morphism $\phi: X \to \P(V) = \P(\Sym^\bullet V^*)$, the projective space of lines in $V$, where we send (at least for $k$-points)
\[P \mapsto (\eval_P:V\to L/\fm_P L\cong k)).\]
For this map $X\to \P(V)$ to make sense, we need $\eval_P$ to always be nonzero, and this happens exactly because $V$ base point free, due to global generation. Also, the isomorphism $L/\fm_P L \cong k$ does not matter, since it would only change $\eval_P$ by scaling.

We should describe this morphism more precisely, i.e., say what happens to non $k$-points and specify the map of structure sheaves. To do this, we will pick coordinates and work locally. Pick a basis $V=\<s_0,\dots,s_r\>$, which in turn gives a basis $V^*=\<s_0^*,\dots,s_r^*\>$. Then, we observe that
\[\eval_P = s_0(P) s_0^* + \cdots s_r(P) s_r^*,\]
so in fact we can describe the morphism by just
\[X \xrightarrow{[s_0: \cdots: s_r]} \Proj k[s_0^*,\dots,s_r^*] \xrightarrow{\cong} \Proj \Sym^\bullet V^*.\]

We then call $L$ \textbf{very ample} if $\phi$ is a closed embedding for some $V$. This won't necessarily be any $V$ that generates $L$, for example take $X=\P^1$, $L=\cO_{\P^1_{[s:t]}}(2)$, and $V=\<s^2,t^2\>$. This $V$ has
\[V\otimes \cO \twoheadrightarrow L\]
because these functions do not simultaneously vanish, so indeed, $V$ induces a morphism $\phi: \P^1 \to \P^1$ by $[s:t]\mapsto [s^2:t^2]$. However, it is \textit{not} a closed embedding, since it's not even injective, e.g., $[-1:1], [1:1]$ map to the same point. On the other hand, if we take $V=\<s^2,st,t^2\>$, then we indeed get a closed embedding, since this is just the Veronese/2-uple embedding of $\P^1 \hookrightarrow \P^2$.

\begin{prop}
    Assume $V\leq H^0(X,L)$ generates $L$, so we it gives a morphism $\phi: X\to \P(V)$. Then, $\phi^* \cO(1) = L$.
\end{prop}
\begin{proof}
    We can see this by choosing a basis $\{s_0,\dots,s_r\}$ of $L$ so $\phi$ is defined as 
    \[X \xrightarrow{[s_0: \cdots : s_r]} \P^r_k.\]
    Then, $\phi^* \cO(1)$ is trivial over each $\phi^{-1}D(x_i)=D(\phi^*x_i)$, since each has nonvanishing section $\phi^*x_i$. In fact, $\phi^{-1}D(x_i)=D(s_i)$. To check this, we show $\phi^{-1}D(x_i) \cap D(s_j) = D(s_is_j)$, so we show
    \[\phi|_{D(s_j)}^{-1}(D(x_i)) = D(s_is_j).\]
    The restricted $\phi$ is a map
    \[\phi|_{D(s_j)}: D(s_j) \to D(x_j) \cong \Spec k[x_{0/j},\dots,x_{r/j}] \hookrightarrow \P^r_k,\]
    defined by $(\frac{s_0}{s_j},\dots,\frac{s_r}{s_0})$, and we have as usual for pulling back functions from affine schemes, $\phi|_{D(s_j)}^{-1}(D(x_{i/j}))=D(\frac{s_i}{s_j}) = D(s_i s_j)$ (where $D(\frac{s_i}{s_j})$ is the usual distinguished open of $\frac{s_i}{s_j}$ inside $D(s_j)$, which then includes into $X$ as $D(s_is_j)$).
    
    So far, we have that both $L,\phi^* \cO(1)$ are trivialized by the same open cover $D(s_i)$, so all that is left is to show that they have the same transition functions. The transition function for $\phi^*\cO(1)$ from $D(s_i)$ to $D(s_j)$ is $\phi^* \frac{x_i}{x_j}$, seen from chasing the trivializations
    \[\cO_{D(s_i)} \xrightarrow{\cdot x_i} \cO(1)_{D(s_i)}, \quad \cO(1)_{D(s_j)} \xrightarrow{\cdot \frac{1}{x_j}} \cO_{D(s_j)}.\]
    Now, $\frac{x_j}{x_i}$ is a function on $D(x_i)$, and it will pullback to a function on $D(s_i)$, where $\phi$ is locally given by the coordinates $(\frac{s_0}{s_i},\dots,\frac{s_r}{s_i})$, so the pull back of the $\frac{x_j}{x_i}$ coordinate is $\frac{s_j}{s_i}$.
    
    On the other hand, the transition function for $L$ from $D(s_i)$ to $D(s_j)$ is, chasing through the trivializations
    \[\cO_{D(s_i)} \xrightarrow{\cdot s_i} L|_{D(s_i)}, \quad L|_{D(s_j)} \xrightarrow{\cdot \frac{1}{s_j}} \cO_{D(s_j)},\]
    given by $\frac{s_i}{s_j}$. Thus, we indeed have $L\cong \phi^* \cO(1)$.
\end{proof}

Now, the closed embedding $X \hookrightarrow \P V = \Proj \Sym^\bullet V^*$ will correspond to some ideal sheaf $\cI_X \leq \cO_{\P V}$, and in turn
\[X \cong \Proj \Sym^\bullet V^*/\Gamma_*(\cI_X),\]
which realizes $X$ as a projective scheme. We study the following commutative diagram.
\begin{cd}
    X \rar["\cong"] \drar[hook] & \Proj \Sym^\bullet V^*/\Gamma_*(\cI_X) \dar[hook]\\
    & \P V
\end{cd}
On one hand, the diagonal map pulls back $\cO_{\P V}(1)$ to $L$. Then, the morphism $\Proj \Sym^\bullet V^*/\Gamma_*(\cI_X) \to \P V$ is induced by the surjective map of graded rings
\[\Sym^\bullet V^* \twoheadrightarrow \Sym^\bullet V^*/\Gamma_*(\cI_X),\]
so therefore it pulls back $\cO_{\P V}(1)$ to $\cO_{\Proj \Sym^\bullet V^*/\Gamma_*(\cI_X)}(1)$. So, in total we have an isomorphism
\[X \cong \Proj S,\]
which pulls back $\cO_{\Proj S}(1)$ to $L$, where $S$ is a finitely generated graded ring over $k$ generated in degree $1$. This statement is in fact an alternative definition for $L$ being very ample.

\begin{prop}
    Assume $X\cong \Proj S$ for $S$ a finitely generated graded ring over $k$ generated in degree $1$. Also assume $\cO_{\Proj S}(1)$ pulls back to $L$. Then, $L$ is very ample. 
\end{prop}

\begin{proof}
    Taking degree 1 generators $x_0,\dots,x_r$ for $S$, we get a surjection of graded rings
    \[k[x_0,\dots,x_r]\twoheadrightarrow S,\]
    which gives a closed embedding
    \[\Proj S \hookrightarrow \P_k^r,\]
    which pulls back $\cO_{\P_k^r}(1)$ to $\cO_{\Proj S}$, so the composition
    \[X \xrightarrow{\cong} \Proj S \to \P_k^r\]
    pulls back $\cO_{\P_k^r}(1)$ to $L$. If we denote the composition $\phi: X\hookrightarrow \P_k^r$, then such a map is given by the section $\phi^* x_0,\dots,\phi^*x_r \in L$, so indeed $L$ is very ample because it gives a map to projective space via its sections.
\end{proof}

\begin{rmk}
    We could have redone this proof using the intrinsic way of embedding $X$ in projective space by instead of surjective $k[x_0,\dots,x_r]\twoheadrightarrow S$, we could have considered $V=S_1^*$, and then considered the surjection
    \[\Sym^\bullet V^* \twoheadrightarrow S.\]
\end{rmk}

Next, we say a projective variety $X\hookrightarrow \P^n$ is \textbf{projectively normal} (equivalently $\cO_X(1)$ is \textbf{normally generated}) if
\[H^0(\P^n, \cO_{\P^n}(k)) \to H^0(X, \cO_X(k))\]
is surjective , i.e., that we don't get any ``extra'' sections not from restriction. We should think of projective normality as saying that resolutions of $\Gamma_* \cO_X$ are the ``same'' as those of $I_X = \Gamma_* \cI_X$. This is in the sense that when building our minimal graded free resolution of $\Gamma_* \cO_X$, we can start by
\[S \twoheadrightarrow \cO_X \to 0,\]
so the next step is generating the kernel, which is the same as just building a resolution of $I_X$. This also says that $\Gamma_* \cO_X = S/I_X$, i.e., just the homogeneous coordinate ring of $X$, and again this need not be true in general.

For an example of a projective variety that is not projectively normal, take the twisted quartic, defined as follows. Take $X=\P^1$, $L=\cO_{\P^1_{[s:t]}}(4)$, and $V=\<s^4,s^3t,st^3,t^4\>\leq H^0(L)$, so $V$ defines
\begin{align*}
    i: \P^1 &\hookrightarrow \P^3\\
    [s:t] &\longmapsto [s^4:s^3t:st^3:t^4],
\end{align*}
which indeed turns out to be a closed embedding. Now, we show that
\[H^0(\P^3, \cO_{\P^3}(1)) \to H^0(X,\cO_X(1))\]
is \textit{not} surjective, by computing dimensions. On one hand, $h^0(\cO_{\P^3}(1))=4$, while
\begin{align*}
    h^0(\cO_X(1)) &= h^0((i_* \cO_{\P^1}) \otimes \cO_{\P^3}(1))\\
    &= h^0(i_*(\cO_{\P^1} \otimes i^* \cO_{\P^3}(1))) \tag{projection formula}\\
    &= h^0(i_*(\cO_{\P^1} \otimes L))\\
    &= h^0(L) = 5,
\end{align*}
and we cannot get a surjection from $\C^4 \to \C^5$. In particular, this tells us $\Gamma_* \cO_X \neq S/I_X$, but instead it is larger, since at least $\cO_X(1)$ has more sections. So, instead the minimal free resolution of $E=\Gamma_* \cO_X$ will be (you can convince yourself?? or use Macaulay2??)
\[0\to S(-3)^3 \to S(-2)^5 \to S(-1) \oplus S \to E\to 0,\]
so while the kernel of $S\to E$ is $I_X$, the map $S\to E$ isn't surjective so we had to do something else, and we're going to get a bigger kernel, and the resolution will not just be that of $I_X$.

\subsubsection{A Characterization of Very Ample}


This is a bit hard to prove---we need some more technology.
\begin{prop}
    Let $k=\overline{k}$ and $X$ a proper $k$-scheme. Let $L$ be a line bundle, and $V\leq H^0(X,L)$ a $k$-subspace such that
    \begin{enumerate}[(1)]
        \item for every pair of distinct closed points $x,y\in X$, there is a section in $V$ which vanishes at $x$ but not at $y$, and
        \item for every closed point $x\in X$ and nonzero tangent vector $\theta \in T_{X/k,x}$ there exists a section $s\in V$ which vanishes at $x$ but whose pullback by $\theta$ is nonzero.
    \end{enumerate}
    Then, $L$ is very ample and the morphism $\phi:X \to \P(V)$ is a closed immersion.
\end{prop}
\begin{proof}
    Condition (1) in particular tells us that $L$ is globally generated, so induces a morphisms $\phi:X\to \P(V)$. Because $X$ is proper/$k$, $\phi$ {\color{red} is proper} (so in particular closed).
    
    Even better, (1) tells us that $\phi$ is an injection on closed points, so $\phi$ will be a homeomorphism onto its image.

    We now need to show
    \[\phi^\sharp: \cO_{\P(V)} \to \phi_* \cO_X\]
    is surjective, which we may check on closed points. Now, $\phi$ injective implies closed points have finite fibers. ``{\color{red} A proper morphism is finite in a neighborhood of a finite fiber}'', so $\phi$ is finite (since $X$ proper).

    % https://stacks.math.columbia.edu/tag/0E8R
    % I think I'm not ready for this yet.
\end{proof}

\begin{cor}
    Let $k=\overline{k}$, $X$ a proper $k$-scheme, and $L$ a line bundle. Assume for every closed subscheme $Z\subseteq X$ of dimension 0 and degree 2 over $k$, the map
    \[H^0(X,L) \to H^0(Z,L|_Z)\]
    is surjective. Then, $L$ is very ample on $X$ over $k$.
\end{cor}

\begin{proof}
    We cite the previous proposition. To get condition (1), we take $Z=x\cup y$, and to get (2), we take, corresponding to a tangent vector $\theta$, $Z$ to be the image of the closed embedding $\theta: \Spec k[\eps]/(\eps^2) \to X$.
\end{proof}

\subsubsection{Products}

Let $L_1,L_2$ be basepoint free line bundles on $X$, giving morphisms
\[\phi_i: X \to \P^{n_i},\]
so we get the product morphism $\phi=(\phi_1,\phi_2): X\to \P^{n_1}\times \P^{n_2}$. Now, we have the commutative diagram
\begin{cd}
    X \rar["\phi"] \ar[rr,bend right,"\phi_i"'] & \P^{n_1} \times \P^{n_2} \rar["\pi_i"] & \P^{n_i}
\end{cd}
which tells us $\phi^* \pi_i^* \cO_{\P^{n_1}}(1) = L_i$. Since pullback commutes with tensors,
\[\phi^*(\cO_{\P^{n_1}} \boxtimes \cO_{\P^{n_2}}) = L_1\otimes L_2.\]
Now, the Segre embedding
\[\P^{n_1} \otimes \P^{n_2} \to \P^{n_1n_2 + n_1 + n_2}\]
is defined by the complete linear series on the line bundle $\cO_{\P^{n_1}}(1)\boxtimes \cO_{\P^{n_2}}(1)$, so the composition
\[X \to \P^{n_1} \otimes \P^{n_2} \to \P^{n_1n_2 + n_1 + n_2}\]
pulls back $\cO_{\P^{n_1n_2 + n_1 + n_2}}(1)$ to $L_1\otimes L_2$.

\subsubsection{Easy Implications of Very Ample}

\begin{prop}
    Let $L$ be very ample on $X$. Then, $L$ is basepoint free. (This is only non-tautological if we take as definition $X\cong \Proj S$ for $S$ f.g. graded $k$-algebra generated in degree 1 with $\cO_{\Proj S}(1)$ pulls back to $L$).
\end{prop}
\begin{proof}
    Since sections of $L$ gives an embedding, in particular they cannot all simultaneously vanish at any point.
\end{proof}

\begin{prop}
    If $L,B$ are line bundles, with $L$ very ample and $B$ basepoint free, then $L\otimes B$ is very ample.
\end{prop}
\begin{proof}
    Say $L$ gives a closed embedding $X\hookrightarrow \P^m$, $B$ gives a morphism $X\to \P^n$. Then, we recall that $L\otimes B$ gives the composite morphism
    \[X \to \P^m \times \P^n \to \P^{mn + m + n},\]
    where the final map is the Segre embedding, which is already a closed embedding. So, we only need to show that $X\to \P^m \times \P^n$ is a closed embedding. Since we know $X\to \P^m \times \P^n \to \P^m$ is a closed embedding, and $\P^m \times \P^n \to (\P^m\times \P^n) \times_{\P^m} (\P^m \times \P^n)$ (the diagonal of the projection) is a closed embedding, by Vakil's cancellation lemma, our desired map is a closed embedding.
\end{proof}

\begin{prop}
    Let $X_1,X_2$ be $k$-schemes, and $L_1,L_2$ very ample. Then, $L_1\boxtimes L_2$ on $X_1 \times X_2$ is very ample. 
\end{prop}

\begin{proof}
    Suppose $L_i$ gives a morphism $\phi_i: X_i \to \P^{n_i}$. We then get a morphism
    \[X_1\times X_2 \to \P^{n_1} \times \P^{n_2}\]
    by $\phi = (\phi_1\circ \pi_1, \phi_2\circ \pi_2)$. By studying the diagram
    \begin{cd}
        X_1 \times X_2 \rar["\phi"] \dar["\pi_i"] & \P^{n_1}\times \P^{n_2}\dar["p_i"]\\
        X_i \rar["\phi_i"] & \P^{n_1}
    \end{cd}
    we see that $\phi^* p_i^* \cO_{\P^{n_i}}(1) = \pi_i^* \phi_i^* \cO_{\P^{n_i}}(1) = \pi_i^* L_i$. So, because pull back commutes with tensors
    \[\phi^*(\cO_{n_1}(1)\boxtimes \cO_{n_2}(1)) = L_1\boxtimes L_2,\]
    and post composing with the Segre embedding
    \[\P^{n_1}\times \P^{n_2} \to \P^{n_1n_2+n_1+n_2}\]
    which pulls back $\cO_{\P^{n_1n_2+n_1+n_2}}(1)$ to $\cO_{\P^{n_1}}(1)\boxtimes \cO_{\P^{n_2}}(1)$, we get that $L_1\boxtimes L_2$ defines the composite
    \[X_1 \times X_2 \to \P^{n_1}\times \P^{n_2} \to \P^{n_1n_2 + n_1 + n_2},\]
    where the last map (the Segre embedding) is already a closed embedding. So, we only need to check that $X_1\times X_2 \to \P^{n_1} \times \P^{n_2}$ is a closed embedding. But this is true, since closed embeddings are ``reasonable'' so are preserved by products.
\end{proof}

\begin{prop}
    Let $X$ be a $k$-scheme and $L$ a very ample line bundle. Then, $L=\cO(D)$ for some effective Cartier divisor $D$.
\end{prop}

\begin{proof}
    Since $L$ is very ample, we get a closed embedding
    \[\phi: X\to \P_k^r\]
    which pulls back $\phi^*\cO(1) = L$. Now, we know that $\cO(1) \cong \cO(H)$ for any hyperplane $H \in H^0(\cO(1))$, since we have
    
    
    {\color{red} TODO}
\end{proof}

\subsection{Many Criteria for Ampleness}

\begin{thm}
    Let $X$ be a proper $k$-scheme, and $L$ a line bundle. The following are equivalent.
    \begin{enumerate}[(a)]
        \item For some $N>0$, $L^{\otimes N}$ is very ample.
        \item[(a')] For all $n\gg 0$, $L^{\otimes n}$ is very ample.
        \item For all finite type quasicoherent sheaves $\cF$, there is $n_0$ such that for $n\geq n_0$, $\cF\otimes L^{\otimes n}$ is globally generated.
        \item As $f$ runs over the global sections of $L^{\otimes n}$ (over all $n>0$), the open subsets $X_f=D(f)$ form a base for the topology of $X$.
        \item[(c')] The open subsets $X_f$ which are affine form a base for the topology of $X$.
        \item[(c'')] The open subsets $X_f$ which are affine cover $X$.
    \end{enumerate}
    In this case, we say $L$ is \textbf{ample} (over $k$).
\end{thm}

\subsubsection{Proof of Criteria}

Some implications are automatic, like (a') $\implies$ (a), (c'') $\implies$ (c') $\implies$ (c).

Assuming (a) and (b) are equivalent, we prove the following, which implies (a) $\implies$ (a').
\begin{prop}
    Let $X$ be a $k$-scheme, and $L,M$ line bundles on $X$. Assume $L$ is ample, then $L^{\otimes n}\otimes M$ is very ample for $n\gg 0$.
\end{prop}
\begin{proof}
    Choose $n_0$ such that for $n\geq n_0$, $M\otimes L^{\otimes n}$ is globally generated. Then, choose $N>0$ such that $L^{\otimes N}$ is very ample. Now, for $n\geq n_0+N$, we get
    \[M\otimes L^{\otimes n} = (M\otimes L^{\otimes(n-N)})\otimes L^{\otimes N}\]
    is the tensor product of a globally generated line bundle and a very ample line bundle, so it is very ample.
\end{proof}


\subsubsection{Some Applications}

\begin{prop}
    Every line bundle on a projective $k$-scheme $X$ is the difference of two very ample line bundles.
\end{prop}

\begin{proof}
    Let $M$ be a line bundle on $X$, and $L$ a very ample line bundle. Then, there exists $n$ such that
    \[M\otimes L^{\otimes n}, \quad L^{\otimes n}\]
    are both very ample, and then
    \[M = (M\otimes L^{\otimes n}) \otimes L^{\otimes (-n)}\]
    is a difference of two very ample line bundles.
\end{proof}

\begin{prop}
    Let $\pi:X\to Y$ a affine morphism of proper $k$-schemes, and $L$ an ample line bundle on $Y$. Then, $\pi^* L$ is ample on $X$.
\end{prop}

\begin{proof}
    We use (c'') to show $\pi^* L$ is ample on $X$. Since $L$ is ample, there exist global sections $\{s_i\}$ where $s_i \in L^{\otimes n_i}$ such that
    \[\bigcup_i Y_{s_i} = Y\]
    and each $Y_{s_i}$ is affine.

    Now, the $\pi^{-1}(Y_{s_i})$ cover $X$, so we want to see that these are affine and complements of the vanishings of sections of $(\pi^* L)^{\otimes n} = \pi^*(L^{\otimes n})$ (since pullback commutes with tensors). Since $\pi$ is finite, it is in particular affine, so $\pi^{-1}(Y_{s_i})$ is affine. Next, as always, we have
    \[\pi^{-1}(Y_{s_i}) = X_{\pi^* s_i},\]
    so these are complements of vanishings of sections.

    So, we indeed have $L$ is ample.
\end{proof}

\begin{prop}
    Let $L$ be ample on a $k$-scheme $X$, and $M$ either ample or basepoint free. Then $L\otimes M$ is ample.
\end{prop}

\begin{proof}
    Choose $N\gg 0$ such that $M^{\otimes N}$ is basepoint free (so if $M$ basepoint free, then $N\geq 1$, and if $M$ is ample, then definitely very ample is enough) and such that $L$ is very ample. Then,
    \[(L\otimes M)^{\otimes N} = L^{\otimes N} \otimes M^{\otimes N}\]
    is the tensor product of very ample with basepoint free, so it is very ample, so $L\otimes M$ is ample.
\end{proof}

\begin{prop}
    Let $X$ be a projective $k$-scheme, $L$ an ample line bundle on $X$, $M$ a line bundle on $X$, and $q_1,\dots,q_m$ points of $X$ (not necessarily closed). Then, for all $N\gg 0$, there is a section of $L^{\otimes N}\otimes M$ that does not vanish on any of the $q_i$.
\end{prop}

\begin{proof}
    First, we argue that we can assume no point is a generalization of any other point. Suppose $p\in \overline{\{q\}}$ (so $q$ is a generalization of $p$), and assume $s\in \Gamma(X,L^{\otimes N}\otimes M)$ is a section which vanishes in the $q$ fiber, and we show $s$ vanishes in $p$ fiber (so the contrapositive says that if we can find $s$ nonvanishing at the specialization, it doesn't vanish in the generalization).

    To see this, we work in a trivialization $\Spec A$ for our line bundle, where $s$ restricts to $a\in A$, and $p,q$ become prime ideals
    \[q \leq p \leq A.\]
    We know $a=0$ in $(A/q)_q$ by assumption, and since $A/q$ is an integral domain, the localization is injective, so $a=0$ in $A/q$. Quotienting further, we get $a=0$ in $A/p$, so $a=0$ in $(A/p)_p$. Thus, $s$ vanishes at $p$.

    Therefore, without loss of generality, we can assume no point is a generalization of any other point.
    
    \textbf{Without Cohomology.} Now, our strategy is to find a sufficiently large $N$ such that there are sections $s_i\in \Gamma(L^{\otimes N}\otimes M)$ which vanish at all points except $q_i$, so that
    \[s = \sum s_i\]
    doesn't vanish at any of the points. We show for $n$ bigger than some $n_0$, there is a section $t=s_1$ which vanishes at $q_2,\dots,q_m$ but not $q_1$. (In general, we pick $n$ to be the max of the $n_0$ for the $q_i$s).
    
    Let $\cI_Z$ be the ideal sheaf of the reduced closed subscheme
    \[Z=\overline{\{q_2\}}\cup\cdots\cup\overline{\{q_m\}}.\]
    By (b), we can choose $n_0$ such that for $n\geq n_0$
    \[L^{\otimes n} \otimes M \otimes \cI_Z \leq L^{\otimes n} \otimes M\]
    is globally generated. Any section of the left hand sheaf, included in $L^{\otimes n} \otimes M$, vanishes in the fiber of $q_2,\dots,q_n$, since in any affine chart $\Spec A$ containing $q_i$ which trivializes $L^{\otimes n} \otimes M$, we get
    \[(L^{\otimes n} \otimes M \otimes \cI_Z)|_{\Spec A} \cong \cI_Z|_{\Spec A} = \tilde{\bigcap_{j=2}^m q_j} \leq \tilde{A},\]
    and so restricting a section $t$ to $\Spec A$, its value in the fiber (considered as a section of $L^{\otimes n}\otimes M$) will be in
    \[0=\left(\left(q_i + \bigcap_{j=2}^m q_j\right)/q_i \right)_{q_i} \leq (A/q_iA)_{q_i} \cong (L^{\otimes n} \otimes M)(q_i),\]
    so indeed $t(q_i)=0$.

    On the other hand, since $q_1 \not\in Z$ (this crucially uses the fact that $q_1$ is not a specialization of the $q_i$), the fiber over $q_1$ is $k\cong (L^{\otimes n}\otimes M)(q_1)$, since
    \[(L^{\otimes n} \otimes M \otimes \cI_Z)|_{X\setminus Z} = (L^{\otimes n} \otimes M)|_Z.\]
    So, by the global generation assumption, we we can pick a global section $t\in \Gamma(L^{\otimes n} \otimes M \otimes \cI)$ which is generates the fiber. Now, indeed, the inclusion of $t$ to a global section of $L^{\otimes n}\otimes M$ is a section which does not vanish at $q_1$ but vanishes at $q_2,\dots,q_m$.
    
    \textbf{Using Cohomology.}  Define $Z$ to be the closed subscheme of $\bigcup \overline{\{q_i\}}$ and $\cI_Z$ the ideal sheaf. Then, study the exact sequence
    \[0\to \cI_Z \to \cO_X \to \cO_Z \to 0,\]
    which we tensor by $L^{\otimes N}\otimes M$ and take cohomology to get
    \[H^0(X,L^{\otimes N}\otimes M) \to H^0(Z,(L^{\otimes N}\otimes M)|_Z) \to H^1(X,L^{\otimes N}\otimes M \otimes \cI_Z).\]
    By Serre vanishing we can get that $H^1$ to vanish by taking $N$ sufficiently large. {\color{red} if the points were all closed points/or maybe all disjoint, then a locally free sheaf on $Z$ would be a direct sum of $\cO_{\mathrm{pt}}$ so surjecting means we can choose whichever values for each point... this argument needs work if they intersect.}
\end{proof}

\begin{prop}
    Suppose $X$ is a projective $k$-scheme, $M$ is a line bundle on $X$, and $q_1,\dots,q_m$ are points of $X$ (not necessarily closed). Show that there are effective Cartier divisors $D_1$ and $D_2$ on $X$ not containing $q_1,\dots,q_m$ such that $M\cong \cO(D_1-D_2)$.
\end{prop}

\begin{proof}
    Since $X$ is projective, we can choose a very ample line bundle $L$, which will then correspond to an effective Cartier divisor.
\end{proof}

\section{Sheaves on Projective Space and Graded Modules}

\subsection{Correspondence}

\subsection{Resolutions}
Given a finitely generated graded module $E$ of dimension 0 over $S=\C[x_0,\dots,x_n]$, by the Auslander-Buchsbaum formula (since it will be finite length, so all $s\in S$ are zero divisors, so regular sequences are length 0, so depth is 0) it has a minimal graded free resolution
\[0\to P_{n+1} \to \cdots \to P_0 \to E \to 0,\]
and the associated sheaf functor $\tilde{(-)}$ is exact, applying it we get an exact sequence
\[0 \to \cP_{n+1} \xrightarrow{d_{n+1}} \cP_n \xrightarrow{d_n} \cdots \to \cP_1 \xrightarrow{d_1} \cP_0 \to 0, \tag{$*$}\]
where $\tilde{E}=0$ because graded dimension 0 (equivalently graded finite length) corresponds to support over the empty set. Notice that we ``forgot'' $E$.

We will study the reverse process, where we start with an exact sequence ($*$) of sheaves on $\P^n$, where each $\cP_i$ is a direct sum of twisting sheaves.

We will show that $H_*^0$ of this sequence is still exact, except for the last term, so there is an extra cokernel, which will be of finite length, recovering $E$. In particular, this will give a correspondence between finitely generated graded $S$-modules of dimension 0 and exact sequences of length $n+1$ on $\P^n$ where matrices are have entries either 0 or nonconstant.

To do this, we split up our LES into SESs as usual
\begin{align*}
    & 0 \to \cP_{n+1} \to \cP_n \to \im d_{n} \to 0 \tag{$\cA_n$}\\
    & 0 \to \im d_n \to \cP_{n-1} \to \im d_{n-1} \to 0 \tag{$\cA_{n-1}$}\\
    & \cdots\\
    & 0 \to \im d_2 \to \cP_1 \to \cP_0 \to 0, \tag{$\cA_1$}
\end{align*}
and we at least get $H_*^0(\cA_k)$ is left exact. We need to show even more that $H_*^0(\cA_k)$ is right exact for $k\geq 2$ (since we only are trying to show exactness everywhere \textit{except} at $\cP_0$).

In other words, we have exact sequences $k=2,\dots,n-1$ (cases $k=2,n$ is slightly different)
\[0\to H_*^0(\im d_{k+1}) \to H_*^0(\cP_k) \to H_*^0(\im d_k) \to H_*^1(\im d_{k+1}),\]
and our goal is to get vanishings of $H_*^1(\im d_{k+1})$, so we can glue these together.

We go through this more carefully here. First, our goal is to show $H_*^1(\im d_3)=0$ to get the exactness of $H_*^0(\cA_2)$. We use the LES of $H_*^\bullet$ on $\cA_3$ giving us
\[H_*^1(\cP_3) \to H_*^1(\im d_3) \to H_*^2(\im d_4),\]
and on $\P^n$ ($n>1$), $H_*^1(\cP_3)=0$ (since it is a direct sum of $\cO(a)$). So, it suffices to prove $H_*^2(\im d_4)=0$.

So, our new goal is to show $H_*^1(\im d_4), H_*^2(\im d_4) = 0$ (the first to get the exactness of $H_*^0(\cA_3)$, the second exactness of $H_*^0(\cA_2)$).

Continuing in a similar fashion, our final goal will be to show
\[H_*^1(\im d_n), H_*^2(\im d_n), \dots, H_*^{n-2}(\im d_n)=0,\]
and we get this by studying the LES of $H_*^\bullet$ on $\cA_n$, which gives
\[H_*^i(\cP_n) \to H_*^i(\im d_n) \to H_*^{i+1}(\cP_{n+1}),\]
where the left and right terms vanish for $i=1,\dots,n-2$ as needed.

So, our conclusion is that we get an exact sequence
\[0 \to H_*^0(\cP_{n+1}) \to H_*^0(\cP_{n}) \to \cdots \to H_*^0(\cP_{1}) \to H_*^0(\cP_{0}) \to \coker H_*^0(d_1) \to 0.\]
Now, if we take the associated sheaves, we get our original exact sequence, where
\[\tilde{\coker H_*^0(d_1)} = \coker \tilde{H_*(d_1)} = 0,\]
(by exactness of $\tilde{(-)}$) which implies that in sufficiently high degrees $\coker H_*^0(d_1)$ is zero, so together with finite generation this gives finite length. \qed

We can refine this argument further---we didn't need intermediate cohomology vanishing of $\cP_1$ or $\cP_0$, so they need not be direct sums of line bundles, any coherent sheaf will do. If we make the length of the exact sequence shorter, i.e.,
\[0\to \cP_r \to \cdots \to \cP_0 \to \cF \to 0\]
for $r<n$ (so the number of maps is 1 smaller than our original setup), we can use intermediate cohomology vanishing of $\cP_0$ to get $H_*^0(\cP_0) \to H_*^0(\cF)$ surjective, so that applying $H_*^0$ we really get an exact sequence, no extra cokernel necessary.

\begin{rmk}
    We still have a finite length ambiguity of sorts
\end{rmk}

% \begin{prop}
%     The minimal graded free resolution of $E$ is of length at most degree $n+1$, with 
% \end{prop}

% \begin{itemize}
%     \item define minimal free resolution of graded $S$-module, betti numbers/table
%     \item Goal: understand betti numbers of a variety embedded in projective space
%     \item give geometry example of points in the plane
%     \item what can we say about canonical curves
% \end{itemize}


\end{document}